% !TEX root = Lecture_Koszul.tex
\clearpage
\section{Homological algebra}\label{sec:minimal-resolutions}

% TODO (ChangeProposal9, Lecture 4 follow-up): After approval, refer to
% def:positive-ideal and eq:degree-zero-projection for A_+ and A/A_+ = S.
% When Ext and Tor are introduced in Lectures 6--7, explain
% Hom_A(M,S) = Hom_S(M/MA_+,S) naturally, and
% Ext_A^i(M,S) = Hom_S(Tor_i^A(M,S),S) naturally (ungraded right modules).
% For a minimal graded resolution, identify the graded Ext groups with the
% S-duals of the generator modules P_i/P_iA_+; check the shift sign carefully.
% This is not higher Hom_S(-,S), which is exact. Specify the Yoneda algebra
% action and its opposite convention when formulating Koszul duality.

\subsection{Free and projective modules}

Until Subsection~\ref{subsec:graded-resolutions}, $A$ is an ordinary $\Bbbk$-algebra and all modules are right $A$-modules.

\begin{definition}\label{def:free-module}
A module $F$ is \defn{free} with basis $\{u_i\mid i\in I\}$ if every $v\in F$ has a unique expression $v=\sum_{i\in I}u_i a_i$, where $a_i\in A$ and only finitely many $a_i$ are non-zero.
Equivalently, $F\cong\bigoplus_{i\in I}A$, with $u_i$ corresponding to $1$ in the $i$th summand.
\end{definition}

A map from a free module is determined by arbitrary images of its basis elements.
In particular, a generating family $\{m_i\mid i\in I\}$ of $M$ gives a surjection
\[
\bigoplus_{i\in I}A\longrightarrow M,\qquad
(a_i)_{i\in I}\longmapsto\sum_{i\in I}m_i a_i.
\]
Every module is therefore a quotient of a free module. If $M$ is finitely generated, the free module can be chosen to be $A^{\oplus m}$ for some integer $m\geq0$.

\begin{definition}\label{def:split-sequence}
A short exact sequence $0\to L\xrightarrow{i}X\xrightarrow{q}Y\to0$ \defn{splits} if there is an $A$-linear map $s:Y\to X$ with $qs=\id_Y$.
The map $s$ is called a \defn{section} of $q$.
\end{definition}

Such a section gives $X=i(L)\oplus s(Y)$: write $x=(x-sq(x))+sq(x)$, and use $qs=\id_Y$ to check that $i(L)\cap s(Y)=0$.
Thus $(l,y)\mapsto i(l)+s(y)$ is an isomorphism $L\oplus Y\to X$ compatible with the inclusion and quotient maps.

\begin{definition}\label{def:projective-module}
A module $P$ is \defn{projective} if, for every surjective $A$-linear map $q:X\to Y$ and every $A$-linear map $f:P\to Y$, there is an $A$-linear map $\widetilde f:P\to X$ with $q\widetilde f=f$ (diagram commutes):
\[
\xymatrix@C=32pt@R=20pt{
& & P\ar[d]^f\ar@{-->}[dl]_{\exists \widetilde f}\\
\text{(exact row:)} & X\ar[r]_{
\forall q} & Y\ar[r] & 0 & .
}
\]
\end{definition}

Free modules have this lifting property: for each basis element $u_i$, choose $x_i\in X$ with $q(x_i)=f(u_i)$, and extend $u_i\mapsto x_i$ $A$-linearly.
Projective modules retain this property without requiring a basis.

For fixed $P$, postcomposition defines the functor $\Hom_A(P,-):\Mod A\to\Mod\Bbbk$, acting on $u:X\to Y$ by
\[
\Hom_A(P,u):\Hom_A(P,X)\longrightarrow\Hom_A(P,Y),
\qquad f\longmapsto u\circ f.
\]

\begin{lemma}\label{lem:Hom-left-exact}
The functor $\Hom_A(P,-)$ is \defn{left exact}, that is, every short exact sequence $0\to L\xrightarrow{i}X\xrightarrow{q}Y\to0$ induces an exact sequence
\[
0\longrightarrow\Hom_A(P,L)\xrightarrow{i\circ-}
\Hom_A(P,X)\xrightarrow{q\circ-}\Hom_A(P,Y).
\]
\end{lemma}

\begin{proof}
Since $i$ is injective, $i\circ f=0$ implies $f=0$.
For $g:P\to X$, the equality $qg=0$ is equivalent to $g(P)\subseteq\Ker q=i(L)$, that is, to a unique factorisation $g=i\widetilde g$ with $\widetilde g:P\to L$ $A$-linear.
\end{proof}

Surjectivity of $q$ does not in general imply surjectivity of $q\circ-$; projectivity is precisely the extra condition.

\begin{proposition}\label{prop:projective-characterisations}
For a right $A$-module $P$, the following are equivalent.
\begin{enumerate}
\item $P$ is projective.
\item The functor $\Hom_A(P,-)$ is exact.
\item Every surjective $A$-linear map $q:X\to P$ has an $A$-linear section.
\item $P$ is isomorphic to a direct summand of a free module.
\end{enumerate}
If $P$ is finitely generated, condition (4) can be replaced by $P\in\add A$.
Direct summands and arbitrary direct sums of projective modules are projective.
\end{proposition}

\begin{supplementary}
\begin{proof}
We prove $(1)\Leftrightarrow(2)$ and $(1)\Rightarrow(3)\Rightarrow(4)\Rightarrow(1)$.

\underline{$(1)\Leftrightarrow(2)$.}
By Lemma~\ref{lem:Hom-left-exact}, exactness is equivalent to surjectivity of
$q\circ-:\Hom_A(P,X)\to\Hom_A(P,Y)$ for every surjection $q:X\to Y$.
This says exactly that every $f:P\to Y$ lifts to $X$.

\underline{$(1)\Rightarrow(3)$.}
Apply the lifting property to $f=\id_P$.

\underline{$(3)\Rightarrow(4)$.}
Choose a free module $F$ and a surjection $q:F\to P$.
A section $s:P\to F$ gives $F=\Ker q\oplus s(P)$.
If $P$ is finitely generated, choose $F=A^{\oplus m}$.

\underline{$(4)\Rightarrow(1)$.}
Identify $P$ with a summand of a free module $F$, with inclusion $j:P\to F$ and projection $r:F\to P$ satisfying $rj=\id_P$.
For $q:X\to Y$ surjective and $f:P\to Y$, freeness gives $h:F\to X$ with $qh=fr$.
Then $\widetilde f=hj$ satisfies $q\widetilde f=frj=f$.

\underline{Summands and sums.}
For a direct summand, use the inclusion and projection as in $(4)\Rightarrow(1)$.
A map from $\bigoplus_{i\in I}P_i$ is determined by its restrictions to the $P_i$; lift each restriction and combine the lifts.
\end{proof}
\end{supplementary}

\begin{example}\label{ex:projective-idempotents}
For an idempotent $e\in A$, $A=eA\oplus(1-e)A$ makes $eA$ projective.
For $A=\Bbbk A_2$, the module $e_{22}A$ is projective of dimension $1$, but not free: non-zero free modules have dimension at least $\dim_\Bbbk A=3$.
\end{example}

The functor $\Hom_A(P,-)$ is covariant in the target.
Resolutions also need the other variable: for fixed $N$, precomposition gives the contravariant functor $\Hom_A(-,N):\Mod A\to\Mod\Bbbk$, $u\mapsto(g\mapsto g\circ u)$.
It reverses arrows, so it sends right exact sequences to left exact ones.

\begin{lemma}\label{lem:Hom-contravariant-left-exact}
Every exact sequence $L\xrightarrow{i}X\xrightarrow{q}Y\to0$ induces an exact sequence
\[
0\longrightarrow\Hom_A(Y,N)\xrightarrow{-\circ q}
\Hom_A(X,N)\xrightarrow{-\circ i}\Hom_A(L,N).
\]
\end{lemma}

\begin{supplementary}
\begin{proof}
Since $q$ is surjective, $g\circ q=0$ implies $g=0$.
For $\psi:X\to N$, the equality $\psi i=0$ is equivalent to $\Im i\subseteq\Ker\psi$.
Since $\Im i=\Ker q$, this is equivalent to a unique factorisation $\psi=g\circ q$ with $g:Y\to N$ $A$-linear.
\end{proof}

In general $-\circ i$ is not surjective, even when $i$ is injective and $q$ is surjective.

\begin{example}\label{ex:Hom-contravariant-not-exact}
Let $A=\Bbbk A_2$, $L_1=e_{11}A/\Bbbk e_{12}$ and $L_2=e_{22}A\cong\Bbbk e_{12}$, and consider the exact sequence $0\to\Bbbk e_{12}\xrightarrow{i}e_{11}A\xrightarrow{q}L_1\to0$.
Apply $\Hom_A(-,L_2)$.
For an $A$-linear map $\varphi:e_{11}A\to L_2$ we have $\varphi(e_{11})=\varphi(e_{11})e_{11}\in L_2e_{11}=0$, so $\varphi=0$ as $e_{11}$ generates $e_{11}A$.
So $\Hom_A(e_{11}A,L_2)=0$, while $\Hom_A(\Bbbk e_{12},L_2)\cong\Bbbk$.
Thus $-\circ i$ is not surjective.
\end{example}
\end{supplementary}

Applying $\Hom_A(-,N)$ to a projective resolution of $M$ and taking cohomology will measure this failure of exactness; this is the subject of Section~\ref{sec:ext-tor}.
If $A$ is finite-dimensional semisimple, every surjection splits by Theorem~\ref{thm:artin-wedderburn}, so every $A$-module is projective.

\subsection{Complexes and resolutions}

\begin{definition}\label{def:chain-complex}\label{def:cochain-complex}\label{not:chain-differentials}
A \defn{chain complex} $C_\bullet$ of right $A$-modules is a sequence of $A$-linear \defn{differentials} $d_h$ with
\[
\cdots\longrightarrow C_{h+1}\xrightarrow{d_{h+1}}C_h
\xrightarrow{d_h}C_{h-1}\longrightarrow\cdots,
\qquad d_h d_{h+1}=0,\quad h\in\Z.
\]
Its \defn{homology} is $H_h(C_\bullet):=\Ker d_h/\Im d_{h+1}$, where $h$ is the \defn{homological degree}.
\end{definition}

Exactness at $C_h$ means $H_h(C_\bullet)=0$.
Writing $C^h:=C_{-h}$ and $\delta^h:=d_{-h}$ turns subscripts into superscripts and reverses the order:
\[
\cdots\longrightarrow C^{h-1}\xrightarrow{\delta^{h-1}}C^h
\xrightarrow{\delta^h}C^{h+1}\longrightarrow\cdots,
\qquad \delta^h\delta^{h-1}=0,
\qquad H^h(C^\bullet):=\Ker\delta^h/\Im\delta^{h-1}=H_{-h}(C_\bullet).
\]
This is a \defn{cochain complex} with \defn{cohomology} $H^h$; algebraically it is the same notion.
A chain complex with $C_h=0$ for $h<0$ is $\cdots\to C_1\to C_0\to0$, and $H_0=C_0/\Im d_1$.

\begin{definition}\label{def:projective-resolution}
A \defn{projective resolution} of $M$ is an exact sequence
\begin{equation}\label{eq:projective-resolution}
\cdots\longrightarrow P_2\xrightarrow{d_2}P_1
\xrightarrow{d_1}P_0\xrightarrow{p_0}M\longrightarrow0
\end{equation}
with every $P_h$ projective.
We write $P_\bullet\to M$; the complex $P_\bullet$ itself ends in $P_0\to0$, and $M$ is not a term of it.
It is \defn{free} if every $P_h$ is free.
\end{definition}

Thus $p_0$ is surjective, $\Ker p_0=\Im d_1$, and $\Ker d_h=\Im d_{h+1}$ for $h\geq1$; equivalently, $H_0(P_\bullet)\cong M$ and $H_h(P_\bullet)=0$ for $h\geq1$.

\begin{proposition}\label{prop:resolutions-exist}
Every right $A$-module has a free resolution, hence a projective resolution.
\end{proposition}

\begin{proof}
For a module $N$, the free module $A^{\oplus N}:=\bigoplus_{n\in N}A$ surjects onto $N$ by $(a_n)\mapsto\sum_{n\in N}na_n$.
Start with $N=M$ and repeat for each successive kernel, composing each surjection with the inclusion of its kernel into the preceding free module; the resulting sequence is exact.
\end{proof}

\subsection{Graded modules, tops and Nakayama's lemma}\label{subsec:graded-resolutions}

From now on $A=\bigoplus_{i\in\Z}A_i$ is locally finite-dimensional non-negatively graded with $S=A_0$ finite-dimensional semisimple, and $A_+=\bigoplus_{i>0}A_i$.
Since $A/A_+\cong S$, we regard $S$ as a graded right $A$-module in degree $0$.
All maps are degree-preserving unless otherwise stated.

\begin{definition}\label{def:graded-projective}
A \defn{graded free module} is a direct sum $\bigoplus_{\lambda\in I}A(-i_\lambda)$ for some set $I$ and $i_\lambda\in\Z$; the generator $1$ of $A(-i)$ has degree $i$.
A graded module is \defn{graded projective} if it has the lifting property of Definition~\ref{def:projective-module} for degree-preserving maps.
\end{definition}

For a graded surjection $q:X\to Y$, an element of $Y_i$ has a preimage in $X_i$, so the lifting proof works for graded free modules; and every graded module is a quotient of a graded free module, by taking homogeneous generators.

\begin{corollary}\label{cor:graded-projective-characterisations}
For a graded right $A$-module $P$, the following are equivalent.
\begin{enumerate}
\item $P$ is graded projective.
\item The functor $\Hom_A^\Z(P,-)$ is exact on graded modules.
\item Every degree-preserving surjection onto $P$ has a degree-preserving section.
\item $P$ is a graded direct summand of a graded free module.
\end{enumerate}
In particular, $Q\otimes_S A(-i)$ is graded projective for every finite-dimensional right $S$-module $Q$ and every $i\in\Z$.
\end{corollary}

\begin{supplementary}
\begin{proof}
The proof of Proposition~\ref{prop:projective-characterisations} applies verbatim with degree-preserving maps.
For the last assertion, Corollary~\ref{cor:semisimple-summands} makes $Q$ a direct summand of $S^{\oplus m}$ for some $m$, so $Q\otimes_S A(-i)$ is a graded direct summand of $A(-i)^{\oplus m}$.
\end{proof}
\end{supplementary}

\begin{definition}\label{not:graded-top}
The \defn{graded top} of a graded right $A$-module $M$ is the graded right $S$-module
\[
\grtop M:=M/MA_+.
\]
\end{definition}

A degree-preserving $f:M\to N$ induces $\grtop f:\grtop M\to\grtop N$, $m+MA_+\mapsto f(m)+NA_+$.

\begin{lemma}\label{lem:graded-radical}
If $A$ is finite-dimensional, then $A_+=\rad A$.
Consequently, $\grtop M=\top M$ for every graded right $A$-module $M$, with $\top M:=M/M\rad A$ carrying the quotient grading.
\end{lemma}

\begin{supplementary}
\begin{proof}
If $A$ is finite-dimensional, then $A_i=0$ for all sufficiently large $i$, so $(A_+)^r=0$ for some integer $r\geq1$.
For any simple right module $L$, the submodule $LA_+$ is either $0$ or $L$.
The latter gives $L=L(A_+)^r=0$, a contradiction.
Thus $A_+$ annihilates every simple module and is contained in $\rad A$.
Since $A/A_+\cong S$ is semisimple, every non-zero element of $S$ acts non-trivially on some simple right $S$-module.
Regarding these as simple $A$-modules shows that $\rad A\subseteq A_+$.
\end{proof}
\end{supplementary}

A graded module $M$ is \defn{bounded below} if $M_{<b}=0$ for some $b$.
Every finitely generated graded module is bounded below: take $b$ to be the least degree among the homogeneous components of a finite generating set.

\begin{lemma}[Graded Nakayama lemma]\label{lem:graded-nakayama}
For a bounded below graded right module $M$,
\[
\grtop M=M/MA_+=0\quad\Longrightarrow\quad M=0.
\]
If $N\subseteq M$ is a graded submodule and $N+MA_+=M$, then $N=M$.
\end{lemma}

\begin{proof}
If $M=MA_+\neq0$, choose the least integer $i$ with $M_i\neq0$.
But $(MA_+)_i=\sum_{r>0}M_{i-r}A_r=0$, a contradiction.
The second assertion follows by applying the first to $M/N$.
\end{proof}

No finite generation is needed: in particular, lifts of homogeneous generators of $\grtop M$ generate $M$.

\begin{supplementary}
\emph{Ungraded analogue.}
The radical $\rad A$ is a two-sided ideal: it consists precisely of the elements acting by zero on every simple right $A$-module.
For a simple module $L$ and $0\neq l\in L$, the map $A\to L$, $a\mapsto la$, is surjective with a maximal right ideal as kernel, so $l\rad A=0$.
Conversely, if $a$ acts by zero on every simple module, apply this to $A/U$ for each maximal right ideal $U$ to get $a\in U$, hence $a\in\rad A$.
Finally $l(ba)=(lb)a=0$ and $l(ab)=(la)b=0$ for $a\in\rad A$, $b\in A$ show $ab,ba\in\rad A$.

\begin{lemma}[Nakayama's lemma]\label{lem:nakayama}
Let $M$ be a finitely generated right $A$-module.
If $M=M\rad A$, then $M=0$.
More generally, if $N\subseteq M$ and $N+M\rad A=M$, then $N=M$.
\end{lemma}

\begin{proof}
For $a\in \rad A$, the right ideal $(1-a)A$ equals $A$: otherwise it lies in a maximal right ideal $U$, and $a,1-a\in U$ gives $1\in U$.
Thus $(1-a)b=1$ for some $b\in A$.

Suppose $M=M\rad A\neq0$ and choose generators $m_1,\ldots,m_r$ with $r$ minimal.
Write $m_r=\sum_{j=1}^r m_j a_j$ with $a_j\in \rad A$ and choose $b$ with $(1-a_r)b=1$.
Then $m_r(1-a_r)=\sum_{j<r}m_j a_j$ and $m_r=\sum_{j<r}m_j a_j b$, contradicting minimality of $r$.
The last assertion follows by applying the first to $M/N$.
\end{proof}

\begin{remark}\label{rem:nakayama-finite-dimensional}
If $A$ is finite-dimensional, then $\rad A$ is \defn{nilpotent}: $(\rad A)^r=0$ for some $r\geq1$.
Indeed, the powers stabilise at $\rad^r A=\rad^{r+1} A$, and Lemma~\ref{lem:nakayama} applied to $(\rad A)^r$ gives $(\rad A)^r=0$.
Thus Nakayama's lemma holds for infinitely generated modules: $M=\rad M$ implies $M=M\rad^r A=0$, and $N+\rad M=M\Rightarrow N=M$ for every submodule $N\subseteq M$.
\end{remark}

For finite-dimensional $A$, every maximal right ideal contains $\rad A$, so $A/\rad A$ is semisimple by Theorem~\ref{thm:artin-wedderburn}.

\begin{definition}\label{def:module-top}
For finite-dimensional $A$, the \defn{top} of a right $A$-module $M$ is $\top M:=M/\rad M$.
\end{definition}

$\top M$ is the maximal semisimple quotient of $M$: it is semisimple, and every map from $M$ to a semisimple $T$ vanishes on $\rad M$ since $\rad T=0$.
A map $f:M\to N$ induces $\top f:\top M\to\top N$, $m+\rad M\mapsto f(m)+\rad N$.
By Nakayama, lifts to $M$ of generators of $\top M$ generate $M$; in particular, $\top f$ surjective implies $f$ surjective when $N$ is finitely generated.

\begin{example}
For $A=\Bbbk A_2$, $\rad A=\Bbbk e_{12}$, $\top(e_{11}A)=e_{11}A/\Bbbk e_{12}=L_1$ and $\top(e_{22}A)=L_2$.
\end{example}
\end{supplementary}

\subsection{Projective covers}\label{subsec:ordinary-covers}

A graded free surjection onto $M$ may use unnecessary generators.
A projective cover removes them.

\begin{definition}\label{def:graded-projective-cover}
A \defn{graded projective cover} of a bounded below lfd graded module $M$ is a surjection $p:P\to M$, with $P$ bounded below, lfd, and graded projective, such that
\[
\grtop p:\grtop P\xrightarrow{\sim}\grtop M.
\]
\end{definition}

Equivalently, $\Ker p\subseteq PA_+$, and no proper graded submodule of $P$ maps onto $M$.

\begin{supplementary}
The identity $p(PA_+)=MA_+$ gives $\Ker(\grtop p)=(\Ker p+PA_+)/PA_+$.
If $\Ker p\subseteq PA_+$ and $p(U)=M$ for a graded submodule $U\subseteq P$, then $U+PA_+=P$; graded Nakayama gives $U=P$.
Conversely, if $\Ker(\grtop p)\neq0$, choose a graded $S$-module complement $W$ to this kernel in $\grtop P$.
The inverse image $U$ of $W$ under $P\to\grtop P$ is a proper graded submodule of $P$.
Since $\grtop p(W)=\grtop M$, we have $p(U)+MA_+=M$; graded Nakayama gives $p(U)=M$.
\end{supplementary}

\begin{proposition}\label{prop:graded-projective-covers}
Every bounded below lfd graded right $A$-module $M$ has a graded projective cover, unique up to graded isomorphism over $M$.
Its domain is isomorphic to
\begin{equation}\label{eq:graded-cover-terms}
P=\bigoplus_{i\in\Z}\bigl((\grtop M)_i\otimes_S A\bigr)(-i),
\end{equation}
where $(\grtop M)_i$ is placed in degree $0$ before shifting.
\end{proposition}

\begin{supplementary}
\begin{proof}
\underline{Existence.}
Since $S$ is semisimple, choose right $S$-linear sections $s_i:(\grtop M)_i\to M_i$ of the quotient maps.
Define on each summand of \eqref{eq:graded-cover-terms}
\[
p:\bigl((\grtop M)_i\otimes_S A\bigr)(-i)\longrightarrow M,
\qquad v\otimes a\longmapsto s_i(v)a.
\]
The map is well defined by $S$-linearity of $s_i$, and preserves degree.
Each summand of $P$ is graded projective, with graded top $(\grtop M)_i$ in degree $i$.
Thus $\grtop p$ identifies $\grtop P$ with $\grtop M$.
Graded Nakayama, applied to $M/p(P)$, makes $p$ surjective.
If $M_i=0$ for $i<b$, then $P_i=0$ for $i<b$, and $P_j=\bigoplus_{b\leq i\leq j}(\grtop M)_i\otimes_S A_{j-i}$ is finite-dimensional.

\underline{Uniqueness.}
For covers $p:P\to M$ and $q:Q\to M$, projectivity of $P$ gives $f:P\to Q$ with $qf=p$.
Since $q(f(P))=M$ and $q$ is a cover, $f(P)=Q$.
Projectivity of $Q$ gives a section $s:Q\to P$ of $f$, and $p(s(Q))=q(Q)=M$ gives $s(Q)=P$.
Thus $f$ and $s$ are inverse isomorphisms.
\end{proof}
\end{supplementary}

\begin{supplementary}
\emph{Ungraded analogue.}

\begin{definition}\label{def:projective-cover}
A \defn{projective cover} of $M$ is a surjective map $p:P\to M$ with $P$ projective such that
\[
U\subseteq P\text{ a submodule and }p(U)=M
\quad\Longrightarrow\quad U=P.
\]
\end{definition}

\begin{proposition}\label{prop:cover-top}
Let $A$ be finite-dimensional, let $P$ be a projective right $A$-module, and let $p:P\to M$ be surjective.
The following are equivalent.
\begin{enumerate}
\item $p$ is a projective cover.
\item $\Ker p\subseteq \rad P$.
\item $\top p:\top P\to\top M$ is an isomorphism.
\end{enumerate}
\end{proposition}

\begin{proof}
\underline{$(1)\Rightarrow(2)$.}
If $\Ker p\not\subseteq\rad P$, there is a maximal submodule $U\subset P$ not containing $\Ker p$.
Then $U+\Ker p=P$, so $p(U)=M$, contradicting (1).

\underline{$(2)\Rightarrow(1)$.}
If $p(U)=M$, then $U+\Ker p=P$, so $U+\rad P=P$ by (2).
Remark~\ref{rem:nakayama-finite-dimensional} gives $P/U=0$.

\underline{$(2)\Leftrightarrow(3)$.}
The equality $p(\rad P)=\rad M$ makes $\top p$ surjective, with kernel $(\Ker p+\rad P)/\rad P$.
\end{proof}

\begin{theorem}\label{thm:ordinary-projective-covers}
Let $A$ be a finite-dimensional algebra.
Every right $A$-module $M$ has a projective cover; if $M$ is finite-dimensional, so is the cover.
Over any algebra, projective covers $p:P\to M$ and $q:Q\to M$ are isomorphic over $M$.
\end{theorem}

\begin{proof}
Uniqueness is the argument of Proposition~\ref{prop:graded-projective-covers}.

\underline{Existence for finite-dimensional $M$.}
Choose a surjection $p:P\to M$ from a finite-dimensional projective module with $\dim_\Bbbk P$ minimal (a quotient of some $A^{\oplus m}$).
Suppose $U\subsetneq P$ and $p(U)=M$.
Projectivity gives $u:P\to P$ with $pu=p$ and $u(P)\subseteq U$.
For $r\geq1$ large, kernels and images of powers of $u$ have stabilised, and $P=\Ker u^r\oplus\Im u^r$.
Indeed, if $u^r(x)\in\Ker u^r$, then $x\in\Ker u^{2r}=\Ker u^r$, so $u^r(x)=0$; equality follows by the dimension formula.
Thus $\Im u^r$ is projective, and $pu^r=p$ gives $p(\Im u^r)=M$.
But $\Im u^r\subseteq U\subsetneq P$ contradicts minimality.

\underline{Existence for arbitrary $M$.}
Write $\top M\cong\bigoplus_{\lambda\in I}L_\lambda$ with simple $L_\lambda$.
Take finite-dimensional covers $q_\lambda:Q_\lambda\to L_\lambda$ and lift each composite $Q_\lambda\to L_\lambda\hookrightarrow\top M$ through $M\to\top M$ to $f_\lambda:Q_\lambda\to M$.
Set $P:=\bigoplus_\lambda Q_\lambda$ and $p:=\sum_\lambda f_\lambda$.
$P$ is projective by Proposition~\ref{prop:projective-characterisations}, and $\top p$ is identified with $\bigoplus_\lambda L_\lambda\xrightarrow{\sim}\top M$.
Hence $M=p(P)+\rad M$, so $p$ is surjective by Remark~\ref{rem:nakayama-finite-dimensional}, and a projective cover by Proposition~\ref{prop:cover-top}.
\end{proof}

\begin{example}\label{ex:ordinary-A2-covers}
For $A=\Bbbk A_2$, the maps $e_{11}A\to L_1$ and $\id:e_{22}A\to L_2$ are projective covers: their kernels are $\rad(e_{11}A)$ and $0$.
The surjection $A=e_{11}A\oplus e_{22}A\to L_1$ killing $e_{22}A$ is not a cover, since its restriction to $e_{11}A$ is already surjective.
\end{example}
\end{supplementary}

\subsection{Minimal resolutions}

\begin{definition}\label{def:minimal-resolution}
A \defn{graded projective resolution} is a projective resolution \eqref{eq:projective-resolution} in which the modules are graded, the differentials are degree-preserving, and the $P_h$ are graded projective.
It is \defn{minimal} if $p_0:P_0\to M$ and the maps $p_h:P_h\to K_h$, $x\mapsto d_h(x)$, with $K_1=\Ker p_0$ and $K_h=\Ker d_{h-1}$ for $h\geq2$, are graded projective covers for all $h\geq1$.
\end{definition}

Homological degree $h$ and grading degree $i$ are distinct: $d_h:(P_h)_i\to(P_{h-1})_i$, and exactness is checked in each grading degree as in Remark~\ref{rem:graded-exactness}.
Since $\Ker p_0=\Im d_1$ and $\Ker d_h=\Im d_{h+1}$, minimality is equivalent to
\begin{equation}\label{eq:graded-minimal-criterion}
\Im d_h\subseteq P_{h-1}A_+\quad\text{for all }h\geq1.
\end{equation}

\begin{theorem}\label{thm:minimal-resolution}
Every bounded below locally finite-dimensional graded right $A$-module $M$ has a minimal graded projective resolution.
Any two such resolutions are isomorphic by degree-preserving isomorphisms $f_h:P_h\to P'_h$ with $p'_0f_0=p_0$ and $d'_hf_h=f_{h-1}d_h$.
\end{theorem}

\begin{proof}
Apply Proposition~\ref{prop:graded-projective-covers} successively to $M$ and to each kernel $K_h$, which remains bounded below and locally finite-dimensional.
For uniqueness, an isomorphism $f_0:P_0\to P'_0$ of the covers of $M$ restricts to an isomorphism $K_1\to K'_1$; uniqueness of covers of the kernels gives $f_1$, and so on.
\end{proof}

By \eqref{eq:graded-cover-terms}, $P_h\cong\bigoplus_{i\in\Z}\bigl((\grtop P_h)_i\otimes_S A\bigr)(-i)$, so $\grtop P_h$ records the types and degrees of the generators needed in homological degree $h$.

\begin{example}\label{ex:minimal-A2}
Let $A=\Bbbk A_2$ with $A_0=S=\Bbbk e_{11}\oplus\Bbbk e_{22}$ and $A_1=\Bbbk e_{12}$, and place $L_1,L_2$ in degree $0$.
Then
\[
0\longrightarrow e_{22}A(-1)\xrightarrow{d_1}e_{11}A
\xrightarrow{p_0}L_1\longrightarrow0,\qquad
d_1(e_{22}a)=e_{12}a
\]
is minimal: $\Ker p_0=\Bbbk e_{12}=(e_{11}A)A_+$ and $d_1$ is an isomorphism onto it, the generator $e_{22}$ and its image $e_{12}$ both having degree $1$.
Its projective terms have tops $L_1$ and $L_2(-1)$.
The other simple module $L_2=e_{22}A$ is already projective.
\end{example}

\begin{example}\label{ex:minimal-polynomial}
Let $A=\Bbbk[x]$ with $\deg x=1$.
The quotient $\Bbbk=A/(x)$ has the minimal graded resolution
\[
0\longrightarrow A(-1)\xrightarrow{d_1}A\xrightarrow{\pi}\Bbbk
\longrightarrow0,\qquad d_1(a)=xa,\quad\pi(f)=f(0).
\]
Multiplication by $x$ is injective with image $xA=\Ker\pi\subseteq AA_+$, and the shift makes it degree-preserving.
The two projective terms have one generator each, in degrees $0$ and $1$.
\end{example}

\begin{supplementary}
Forgetting the grading need not preserve projective covers.
In Example~\ref{ex:minimal-polynomial}, the proper ungraded submodule $(x-1)A\subset A$ still maps onto $\Bbbk$ under $\pi$, since $\pi(x-1)=-1$.
Thus $\pi$ is a graded projective cover but not an ordinary projective cover.
\end{supplementary}

\begin{example}\label{ex:minimal-exterior}
Give $E=\Bbbk[\epsilon]/(\epsilon^2)$ the grading $\deg\epsilon=1$.
Then
\[
\cdots\longrightarrow E(-3)\xrightarrow{d_3}E(-2)
\xrightarrow{d_2}E(-1)\xrightarrow{d_1}E\xrightarrow{\pi}\Bbbk
\longrightarrow0,
\qquad d_h(a)=\epsilon a,
\]
is a minimal graded resolution: multiplication by $\epsilon$ has kernel and image both equal to $\Bbbk\epsilon$, so it is exact, and $\Im d_h=\epsilon E(-h+1)\subseteq E(-h+1)E_+$.
For every $h\geq1$, the generator of $E(-h)$ and its image under $d_h$ both have degree $h$.
The term in homological degree $h$ has one generator, in degree $h$, and the resolution never terminates.
\end{example}

\begin{advanced}
\begin{remark}\label{not:graded-betti}
The multiplicities in the graded tops are sometimes called \defn{graded Betti numbers}.
For representatives $L_1,\ldots,L_n$ of the simple right $S$-modules, one can write
\[
(\grtop P_h)_i\cong\bigoplus_{y=1}^nL_y^{\oplus\beta_{h,i,y}(M)},
\qquad
\beta_{h,i}(M):=(\beta_{h,i,y}(M))_{1\leq y\leq n}.
\]
Here $h\geq0$ is the homological degree and $i\in\Z$ is the grading degree.
These multiplicities are independent of the chosen minimal resolution by Theorem~\ref{thm:minimal-resolution}.
We will use the modules $\grtop P_h$ directly, without this additional indexing.
\end{remark}
\end{advanced}

\begin{supplementary}
\emph{Ungraded analogue.}
For a finite-dimensional algebra $A$ and a right $A$-module $M$, replace graded covers by projective covers in Definition~\ref{def:minimal-resolution}.

\begin{definition}\label{def:ordinary-minimal-resolution}
A projective resolution \eqref{eq:projective-resolution} is \defn{minimal} if $p_0:P_0\to M$ and the surjections $p_h:P_h\to K_h$, $x\mapsto d_h(x)$ (with $K_1=\Ker p_0$, $K_h=\Ker d_{h-1}$ for $h\geq2$) are projective covers.
\end{definition}

Each differential then factors as a projective cover onto its image followed by an inclusion:
\[
\xymatrix@C=30pt@R=30pt{
\cdots\ar[r] & P_2\ar[rr]^{d_2}\ar@{->>}[dr]|{\substack{p_2\\\text{proj.\ cover}}}
& & P_1\ar[rr]^{d_1}\ar@{->>}[dr]|{\substack{p_1\\\text{proj.\ cover}}}
& & P_0\ar[rr]|{p_0}_{\text{proj.\ cover}} && M\ar[r] & 0\\
& & K_2\ar@{^{(}->}[ur] && K_1\ar@{^{(}->}[ur]
}
\]
By Proposition~\ref{prop:cover-top}, minimality is equivalent to
\begin{equation}\label{eq:ordinary-minimal-criterion}
\Im d_h\subseteq \rad P_{h-1}\quad\text{for all }h\geq1.
\end{equation}

\begin{theorem}\label{thm:ordinary-minimal-resolutions}
Every right module over a finite-dimensional algebra has a minimal projective resolution; its terms are finite-dimensional if $M$ is.
Over any algebra, two minimal resolutions $P_\bullet\to M$ and $P'_\bullet\to M$ are isomorphic: there are isomorphisms $f_h:P_h\to P'_h$ with $p'_0f_0=p_0$ and $d'_h f_h=f_{h-1}d_h$.
\end{theorem}

\begin{proof}
As in Theorem~\ref{thm:minimal-resolution}, using Theorem~\ref{thm:ordinary-projective-covers}; if $M$ is finite-dimensional, choose finite-dimensional covers, whose kernels are again finite-dimensional.
\end{proof}

More generally, maps $f_h$ with these identities form a \defn{map of resolutions}, that is, a commutative diagram
\[
\xymatrix@C=24pt@R=22pt{
\cdots\ar[r] & P_1\ar[r]^{d_1}\ar[d]^{f_1}
& P_0\ar[r]^{p_0}\ar[d]^{f_0} & M\ar[r]\ar[d]^{\id_M} & 0\\
\cdots\ar[r] & P'_1\ar[r]^{d'_1}
& P'_0\ar[r]^{p'_0} & M\ar[r] & 0.
}
\]
The kernels $\Ker p_0$ and $\Ker d_h$ are the successive \defn{syzygies} of $M$; their isomorphism classes do not depend on the minimal resolution.
Over an arbitrary algebra, Proposition~\ref{prop:resolutions-exist} still gives a free resolution, but projective covers of the kernels need not exist; the obstruction is to existence of minimal resolutions, not to uniqueness.

\begin{example}\label{ex:ordinary-A2-resolution}
For $A=\Bbbk A_2$ the covers in Example~\ref{ex:ordinary-A2-covers} give the minimal resolution of $L_1$,
$0\to e_{22}A\xrightarrow{d_1}e_{11}A\xrightarrow{p_0}L_1\to0$ with $d_1(e_{22}a)=e_{12}a$; this is Example~\ref{ex:minimal-A2} without the shift.
\end{example}

\begin{example}\label{ex:ordinary-exterior-resolution}
Let $E=\Bbbk[\epsilon]/(\epsilon^2)$.
Every $a+b\epsilon$ with $a\neq0$ has inverse $a^{-1}-a^{-2}b\epsilon$, so $\rad E=\Bbbk\epsilon$.
The sequence $\cdots\xrightarrow{\epsilon\cdot-}E\xrightarrow{\epsilon\cdot-}E\to\Bbbk\to0$ is a minimal resolution of $\Bbbk$ with all terms non-zero: a finite-dimensional algebra can have a simple module with an infinite minimal resolution.
\end{example}
\end{supplementary}

\begin{definition}\label{def:noetherian-algebra}
$A$ is \defn{right noetherian} if every submodule of every finitely generated right $A$-module is finitely generated; \defn{left noetherian} is defined symmetrically, and \defn{noetherian} means both.
\end{definition}

\begin{supplementary}
\begin{remark}\label{rem:noetherian-acc}
$A$ is right noetherian if and only if every ascending chain $I_1\subseteq I_2\subseteq\cdots$ of right ideals stabilises (the \defn{ascending chain condition}, ACC); similarly on the left.
Finite-dimensional algebras and $\Bbbk[x_1,\ldots,x_m]$ are noetherian, whereas the free algebra $\Bbbk\langle x,y\rangle$ is not.
Thus the hypothesis of Corollary~\ref{cor:noetherian-minimal-resolution} is not automatic.

\emph{Proof.}
If submodules of finitely generated modules are finitely generated, the union of a chain $I_1\subseteq I_2\subseteq\cdots$ of right ideals is finitely generated, with generators in some $I_k$; so the chain stabilises at $I_k$.
Conversely, assume the ACC.
Every right ideal $I$ is finitely generated: otherwise choose inductively $u_{k+1}\in I\setminus(u_1A+\cdots+u_kA)$, giving a strictly ascending chain.
We show by induction on $m$ that every submodule $U\subseteq A^{\oplus m}$ is finitely generated; the case $m=1$ was just proved.
For $m>1$, let $\pi:A^{\oplus m}\to A$ be the projection onto the last coordinate.
Then $U\cap\Ker\pi\subseteq A^{\oplus(m-1)}$ and $\pi(U)$ are finitely generated, and generators of $U\cap\Ker\pi$ together with lifts to $U$ of generators of $\pi(U)$ generate $U$.
Finally, a finitely generated module is a quotient $A^{\oplus m}/K$, and each submodule of it has the form $U/K$ with $K\subseteq U\subseteq A^{\oplus m}$, hence is finitely generated.
\end{remark}
\end{supplementary}

\begin{corollary}\label{cor:noetherian-minimal-resolution}
Assume that $A$ is right noetherian.
Every finitely generated graded right $A$-module has a minimal graded projective resolution whose terms and successive kernels are finitely generated.
In particular, this applies to $S=A/A_+$.
\end{corollary}

\begin{proof}
A finitely generated graded module is bounded below and lfd, being a quotient of a finite direct sum of shifts of $A$.
Its $\grtop M$ is finite-dimensional, so its graded projective cover is finitely generated by \eqref{eq:graded-cover-terms}.
Right noetherianity makes the kernel finitely generated, so the same argument applies to every successive kernel.
\end{proof}

Without noetherianity, boundedness below and local finite-dimensionality still give existence of a minimal graded resolution (Theorem~\ref{thm:minimal-resolution}), but its terms need not be finitely generated.

\subsection{Exercises}

\begin{exercise}
Let $A=\Bbbk[x]$ with $\deg x=1$, and let $M=\Bbbk[x,x^{-1}]$ be the graded $A$-module with $M_i=\Bbbk x^i$ for all $i\in\Z$ and the action by multiplication.
Show that $M=MA_+\neq0$, so $\grtop M=0$.
Which hypothesis of Lemma~\ref{lem:graded-nakayama} fails?
\end{exercise}

\begin{exercise}
Let $A=\Bbbk[u,v]/(uv)$ with $\deg u=\deg v=1$ and $S=A_0=\Bbbk$.
Show that the minimal graded projective resolution of $\Bbbk=A/A_+$ has $P_0=A$ and $P_h\cong A(-h)^{\oplus2}$ for all $h\geq1$.

\emph{Hint.} Show that $ua+vb=0$ in $A$ if and only if $a\in vA$ and $b\in uA$.
\end{exercise}

\begin{exercise}
Let $A=\Bbbk A_3$, graded by path length.
Find the minimal graded projective resolutions of $L_1$, $L_2$ and $L_3$ and the tops of their terms.
\end{exercise}

\begin{exercise}
Let $A=\Bbbk[x]$ with $\deg x=1$.
\begin{enumerate}
\item Prove that $\rad A=0$ whereas $A_+=(x)$; in particular, Lemma~\ref{lem:graded-radical} fails for infinite-dimensional $A$.
\item Prove that the ungraded module $\Bbbk=A/(x)$ has no projective cover.
\end{enumerate}
\end{exercise}

\begin{hidden}
\emph{Solution to the last exercise.}
(1) If $0\neq f\in\rad A$, then $xf\in\rad A$, so $1-xf$ has a right inverse.
This is impossible because $1-xf$ has positive degree.
(2) Suppose $p:P\to\Bbbk$ is a projective cover, and choose $u\in P$ with $p(u)=1$.
Then $p(uA)=\Bbbk$, so $P=uA$.
The surjection $A\to P$, $a\mapsto ua$, splits because $P$ is projective.
The resulting projection $A\to A$ is multiplication by an idempotent $e\neq0$, and $P\cong eA$.
Since $\Bbbk[x]$ has no idempotents other than $0$ and $1$, we have $P\cong A$.
But every surjection $A\to\Bbbk$ maps the proper submodule $(x-1)A$ onto $\Bbbk$, contradicting the cover property.
\end{hidden}
